\section{Guiding Principles}\label{sec:principles}

This draft treats pedagogical, governance, and operational concerns together, following the three dimensions of the AI ecological education policy framework proposed by \citet{chan2023}, and is guided by the national principles in \cmo{} \citep{cmo21}. The following principles guide interpretation of every provision.

\begin{enumerate}
  \item \provision{Human-centered and accountable.} People, not AI systems, remain responsible for decisions, outputs, grades, and published work. AI supports professional judgment; it does not replace it.
  \item \provision{Enabling by default.} Lawful and ethical AI use consistent with course or work objectives is permitted and encouraged without prior approval, unless a higher-risk tier applies.
  \item \provision{Proportionality.} Requirements increase only as risk increases (to personal data, assessment validity, or consequential decisions). Low-risk uses carry minimal obligations.
  \item \provision{Transparency through simplicity.} Disclosure is expected where AI use affects assessment or authorship, using short standard statements rather than lengthy reports.
  \item \provision{Learning first.} Students should still develop the knowledge and skills each course intends. AI use is judged by whether it supports or undermines the learning outcomes.
  \item \provision{Integrity through design.} Clear expectations, authentic assessment, and process evidence are preferred to surveillance and detector-driven enforcement; \cmo{} provides that detection tools shall not be the sole basis for determining academic misconduct (Sec.~9).
  \item \provision{Fairness and inclusion.} Policies must not disadvantage students by language, disability, connectivity, or ability to pay for tools.
  \item \provision{Privacy and security.} Personal and confidential data are protected in accordance with law and University policy.
  \item \provision{Academic freedom and professional growth.} Faculty exercise professional discretion in course design and are supported to experiment and to learn.
  \item \provision{Continuous learning.} AI capabilities and risks change quickly; the framework is reviewed regularly and revised through evidence and feedback.
\end{enumerate}
